% Artículo VII. Incorporación de resultados preexistentes, 10-09-2026.
% Propietarios íntegros preservados en fuentes_conservadas:
% 17_pantalla_cubica_soldadura.tex (P), 18_clausura_energetica.tex (E).
% P: incidencia, cociente, soldadura celular y refinamiento.
% E: respuesta, unicidad, extensión mínima y expectativa tracial.
% Dependencias anteriores de VII: núcleo APP--TRIT--TPK; realización
% cúbica de la ruta, cociclo de transporte y compresión q_vis=5/9.
% Este fragmento no constituye por sí solo el artículo autónomo.
% La identidad de pantalla se distingue de la contracción material de espín.

\subsection{Incidence cubique et soudure de frontière}
\label{vii:sec:pantalla-respuesta}

La réalisation cubique du transport conserve six incidences orientées.
Son quotient géométrique, la forme lorentzienne et la réponse énergétique
s’obtiennent à partir de la même matrice d’incidence. Le transport provient du
chemin enrichi APP--TRIT--TPK : la face, le signe, la retenue et la
mémoire accompagnent chaque arête. Cette section développe
l’opération linéaire qui agit sur cette représentation de frontière.

\subsubsection{Incidence cubique et quotient de rang quatre}
\label{vii:sec:cociente-pantalla}

Soit \(H_F=\mathbb R^F\), avec
\(F=\{(i,\epsilon):1\leq i\leq3,\ \epsilon\in\{-1,1\}\}\),
son produit scalaire euclidien et sa base \(e_{i,\epsilon}\).
L'opposition des faces est l'involution orthogonale
\(Re_{i,\epsilon}=e_{i,-\epsilon}\). Nous définissons
\[
 u=\sum_{i,\epsilon}e_{i,\epsilon},\qquad
 P_u=\frac{uu^{\mathsf T}}6,\qquad P_-=\frac{I-R}{2},\qquad
 P_0=\frac{I+R}{2}-P_u,\qquad P_\square=P_u+P_-.
\]
Les sommets sont \(V=\{-1,1\}^3\). La matrice d'incidence \(M:\mathbb R^V\longrightarrow H_F\) est déterminée par
\[
 M_{(i,\epsilon),\nu}=\mathbf1_{\{\nu_i=\epsilon\}}.
\]

\begin{proposition}[Décomposition de l'incidence]
\label{vii:prop:incidencia}
On a
\[
 MM^{\mathsf T}=12P_u+4P_-,
 \qquad
 \ker M^{\mathsf T}=P_0H_F.
\]
Par conséquent, le quotient
\(Q_\square=Q=H_F/P_0H_F\), identifié orthogonalement à
\(P_\square H_F\), est de dimension quatre.
\end{proposition}

\begin{proof}
Une face contient quatre sommets. Deux faces opposées n'ont pas de sommets
communs et deux faces d'axes différents ont deux sommets communs.
En écrivant \(\mathbf J=uu^{\mathsf T}\), on obtient
\[
 MM^{\mathsf T}=4I+2(\mathbf J-I-R)
 =2(I-R)+2\mathbf J=4P_-+12P_u.
\]
Les trois projecteurs \(P_u,P_-,P_0\) sont orthogonaux et leur somme vaut \(I\) ; leurs rangs sont respectivement un, trois et deux. L'identité \(\|M^{\mathsf T}x\|^2=\langle x,MM^{\mathsf T}x\rangle\) identifie le noyau annoncé. Le rang complémentaire est quatre.
\end{proof}

Sur \(Q\), la forme
\[
 B=P_--P_u
\]
a pour signature \((3,1)\). L'orientation temporelle est fixée par \(u\).
Avec
\[
 t=\frac{u}{\sqrt6},\qquad
 d_i=\frac{e_{i,+}-e_{i,-}}{\sqrt2},\qquad
 W=\sqrt6P_u+\sqrt2P_-,
\]
on a \(We_{i,\epsilon}=t+\epsilon d_i\). Ces six vecteurs
sont nuls pour \(B\), tandis que
\[
 3t+\nu_1d_1+\nu_2d_2+\nu_3d_3
\]
a pour norme quadratique \(-6\). La même incidence satisfait
\(MM^{\mathsf T}=2W^*W\) sur \(Q\). Les rotations propres du cube
commutent avec les projecteurs, préservent les deux formes et fixent \(t\).
Ces identités déterminent la métrique de la représentation cubique.

\subsubsection{Soudure cellulaire et transport de mémoire}
\label{vii:sec:soldadura-celular}

Pour chaque arête orientée \(a\), la lecture du chemin fournit
la face \(\lambda(\underline a)\), le signe
\(\sigma(a)\), les fibres d’origine et de destination et le transport
\(U(a):Q_{s(a)}\longrightarrow Q_{t(a)}\).
L’inversion vérifie \(U(\bar a)=U(a)^{-1}\).
Les demi-arêtes de frontière conservent leur inversion formelle et leur incidence
sur le bord. Avec \(\ell_m=\ell_0\,9^{-m}>0\), la soudure est
\begin{equation}
 \beta_m(a)=\ell_m\sigma_m(a)
              n_{\lambda(\underline a)}^{s(a)},\qquad
 n_{i,\epsilon}=t+\epsilon d_i.
 \label{vii:eq:soldadura-arista}
\end{equation}
La face et le signe sont transportés avec la mémoire de la route. La convention d'inversion exige
\begin{equation}
 \beta_m(\bar a)=-U_m(a)\beta_m(a).
 \label{vii:eq:inversion-soldadura}
\end{equation}
En appliquant deux fois cette relation, on retrouve \(\beta_m(a)\) ; le registre de retenue conserve séparément la composition ordonnée.

\begin{proposition}[Naturalité de la soudure par raffinement]
\label{vii:prop:refinamiento-soldadura}
Soit \(a=a_1\cdots a_9\) un raffinement compatible d’une arête.
Supposons que les neuf incidences filles, transportées vers la fibre du
prédécesseur, conservent la face et le signe, et que
\(\ell_{m+1}=\ell_m/9\). Soit
\(\mathcal U_j:Q_{m,s(a)}\longrightarrow Q_{m+1,s(a_j)}\)
le transport cumulé qui inclut l’identification par raffinement
et le trajet jusqu’à l’origine du successeur \(j\). Alors
\[
 \beta_m(a)=
 \sum_{j=1}^{9}\mathcal U_j^{-1}\beta_{m+1}(a_j).
\]
L'égalité est compatible avec l'inversion et avec la concaténation du registre de mémoire.
\end{proposition}

\begin{proof}
Chaque terme, exprimé dans la fibre de l’origine du prédécesseur, est
\((\ell_m/9)\sigma_m(a)n_{\lambda(\underline a)}\).
La somme de neuf termes reproduit la soudure du prédécesseur.
L’inversion renverse l’ordre des transports et remplace chacun
par son inverse ; la relation \eqref{vii:eq:inversion-soldadura} fournit
le signe global. La mémoire se concatène dans ce même ordre, de sorte
que le raffinement conserve tant la représentation vectorielle que son
registre enrichi.
\end{proof}
